% before: ib-host-group
The window bound is proved inside a host field whose Galois group is
$\{\pm1\}^2\times A$ for a finite abelian group $A$. The balance identity uses
only the group structure, the index-two subgroup fixing $\sqrt{-d}$, complex
conjugation and the inertia at $d$, so we state it for an abstract group of
this shape.

% before: ib-host-group-fields
The abstract elements of the host group have arithmetic meaning in a host field.
The dictionary explains why primes whose Frobenius lies outside the split
subgroup are called inert.

% before: ib-arithmetic-data
The identity needs, for each prime, a Frobenius class and an inertia subgroup,
and only a few facts about them: the primes in a finite set $S$ split, the
primes outside $S\cup\{d\}$ are unramified, and $d$ is tamely ramified with
inertia of order two. These are the facts that orthogonality of characters
turns into the shape of the inert-prime sum.

% before: ib-host-chebyshev
Each character of the host group has a corrected Chebyshev sum, built from its values at
prime powers and corrected by a number $\ell(\psi)$. For a host field and the
right choice of $\ell$ these are the corrected Chebyshev sums of
Chapter~\ref{chapter-windowed-landau}, but the identity holds for every
$\ell$.

% before: ib-inert-sums
The identity compares three quantities attached to a scale $a$, a class $t$
and a weight function $w$ on the split subgroup. The inert-prime sum is the
quantity that divisibility of periods bounds from below, and the main term is
the one that covolumes bound from above.

% before: ib-prime-power-orthogonality
The values of characters at prime powers are averages over cosets of the
inertia subgroup. Orthogonality therefore picks out the coset of the Frobenius
class, with a factor $2$ from the index of the split subgroup.

% before: ib-inert-expansion
The first half of the identity writes the inert-prime sum as a sum over all
characters. Each prime contributes according to its type: split primes
contribute nothing because their Frobenius classes and inertia lie in the split
subgroup, $d$ contributes once per power, and an unramified prime contributes
only when it is inert.

% before: ib-truncation
The second half rescales each character sum to the corrected Chebyshev sum at the scale
$a/W$. The rescaling produces an explicit remainder, which for the trivial
character carries the archimedean terms $a\log a-(1+\gamma)a$.

% before: ib-balance-identity
Combining the two halves, the remainders add up to the main term, because
summing a character over the complement of the split subgroup isolates the
trivial character and $\chi_C$. The identity is exact, with no error term.

% before: ib-comparison-data
The inequality that bounds the inert-prime sum is stated for numbers, not for
lattices: slopes, heights, sizes and losses attached to base points $u_0$ and
words $u$. Chapters~\ref{chapter-hermitian-slopes} to~\ref{chapter-covolumes}
construct the lattices and prove the hypotheses for their invariants.

% before: ib-comparison-quantities
The inequality and its hypotheses are expressed through a few derived
quantities. Two of them carry the argument: the weighted mean slope $s(u_0)$
of the words, and the main term $\mathrm{Main}(u_0)$, which in the application
is the sum of the main terms of the balance identity over the top classes and
the lifts.

% before: ib-comparison-hypotheses
There are eleven hypotheses. The division inequality (H2) and the frame slope
inequality (H9) are the two measurements of the glued lattice; (H7) and (H8)
bring in the gain from divisibility; the others bound the losses, heights and
sizes that enter the error.

% before: ib-slope-of-n-lower
The first step bounds the mean slope of the words from below. Divisibility
raises each slope by the gain, and the division inequality converts the height
of each word into a lower bound in terms of the slope of the glued lattice.

% before: ib-spread
The frame slope inequality loses a term that depends on how far the largest
slope exceeds the mean slope. The hypotheses on the words bound this spread by
the variation of their exponents, heights and sizes.

% before: ib-comparison-inequality
Putting the two bounds for the slope of the glued lattice together, the slope
cancels and the gain is bounded by the main term. Since the gain bounds the
inert-prime sum up to a loss, the inert-prime sum is bounded by the main term
plus an explicit error.

% before: ib-two-measurements
The comparison inequality is the descendant of Stepanov's method in this
argument. The remark identifies which hypotheses measure the glued lattice from
below and which from above.

% before: ib-size-windows
To turn the balance identity into a bound for a weighted Chebyshev sum, the
weights $w$ must be close to a fixed size $b$ at the thick places and small
elsewhere, and all corrected Chebyshev sums must be small at very large
scales.
These size conditions make the error terms of the next lemmas smaller than the
main quantities.

% before: ib-kept-characters
Only the characters trivial on a subgroup $A_0$ survive averaging over $A_0$.
Among them, the kept characters are those whose corrected Chebyshev sums enter the
weighted sum $G_K$, and the conditions say that the top classes $T$ are
equidistributed for them and that their values $\ell(\psi)$ are bounded.

% before: ib-weight-lipschitz
Replacing a weight $w$ near $b$ by $b$ itself changes $w\,g_\psi(A/w)$ by an
amount proportional to $|w-b|$. The constant comes from Chebyshev's bound for
$\sum_{n\leq Y}\Lambda(n)$.

% before: ib-chebyshev-lipschitz
The scales $a(t)/b$ differ from $x$ by at most $1$. Moving the scale from
$a(t)/b$ to $x$ changes each corrected Chebyshev sum by a logarithmic amount, by Mertens'
estimate.

% before: ib-thick-replacement
After the weights are replaced, the character side of the balance identity
becomes a sum over the thick places at the single scale $A/b$. The other
elements of $\Psi_H$ contribute little because their scales $A/w$ are large, where the
corrected Chebyshev sums are small.

% before: ib-kept-character-sum
Summed over the top classes, each character trivial on $A_0$ contributes
either its weighted corrected Chebyshev sum at $x$, if it is kept, or at least minus its
absolute value. The error per character is the cost of moving the scale and of
the imperfect equidistribution of $T$.

% before: ib-size-arithmetic
The averaged inequality has the form "density times weighted sum $\leq$ error".
The size conditions make the error a factor $x^{1/2+O(\varepsilon)}$ times the
density, which is the arithmetic of this lemma.

% before: ib-core-window-bound
The four previous lemmas combine into a bound for the weighted sum
$g^{\mathcal K}_c$ of the corrected Chebyshev sums of the kept characters. Its
hypothesis is that the comparison inequality holds at
every class of $A_0$, in the form of a bound for the gap between the inert-prime
sums and their main terms.

% before: ib-window-bound-shape
The proof hides the shape of the argument in its error terms. The remark states
the approximate identity that the proof makes rigorous.

% before: ib-host-arithmetic-data
In a host field the Frobenius and inertia data come from the Artin map on
$(\mathbf Z/Q_Sqd\mathbf Z)^\times$. The primes dividing $Q_Sq$ split in
$\mathbf Q(\sqrt{-d})$, and $d$ is ramified with inertia generated by $\sigma$,
which are the conditions of the abstract arithmetic data.

% before: ib-host-character-values
The corrected Chebyshev sums of the host group must be identified with those of
Dirichlet characters. Each character of the Galois group determines a primitive
Dirichlet character through the Artin map, and the compatibility of $\ell$ makes
the two sums agree.

% before: ib-level-sum-host-characters
The weighted sum of the kept host characters must be the weighted Chebyshev sum
$G_K$ of Chapter~\ref{chapter-windowed-landau}, with the same weights. The host
conditions match the kept characters with $K$, and the weights agree because a
kept character trivial on $A$ is the trivial character or the one attached to
the exceptional character.

% before: ib-host-kept-characters
The remaining conditions on the kept characters, and the bound on the others,
are what the host conditions of Chapter~\ref{chapter-host-fields} were chosen
to give. The characters that are not kept have square-root cancellation at
the scale $x$.

% before: ib-error-windows
The error of the comparison inequality is a long expression. Under size
conditions on the parameters, which the construction satisfies at every large
scale, it is bounded by a simple expression.

% before: ib-error-size
With these conditions the error is at most a power $x^{O(\varepsilon)}$ times a
quantity linear in the parameters. This puts it in the form of the error
bracket of the size windows.

% before: ib-gap-bound
When the numbers of the comparison data are those of the balance identity, the
quantity $\mathrm{ASK}-\mathrm{Main}$ of the comparison inequality is the gap of
Definition~\ref{ib-kept-characters}. The comparison inequality and the bound on
its error then give the gap bound required by
Theorem~\ref{ib-core-window-bound}.

% before: ib-core-data
All the data needed at one host field and one scale fit into a single
definition. A family of definite-side data that is tied at a scale carries such
core data there, and Chapter~\ref{chapter-assembly} shows that suitably designed
families are tied at all large scales.

% before: ib-host-window-bound
At one host field and one scale, the host conditions and the core data now give
the window bound. The exponent $e_D+e_L+7$ collects the exponents of the
division and loss bounds and the losses of the averaging.

% before: ib-window-bound-on-scales
The window bound along the scales of a family is condition (1) of
Theorem~\ref{wl-level-from-windows}. It follows from the bound at one host field,
because a family tied at a scale carries core data there.
